\documentclass{article}
\usepackage{pathfinder-readable}

\title{Can Pairwise Judgements Find Better Quotes in a Double Auction?}

\begin{document}
\maketitle

\begin{abstract}
One paper studies language-model traders in a double auction; the other uses pairwise judgements to assign
credit to cooperating agents. The thread asked whether an adapted pairwise judge could identify quotes that
improve the auction's final welfare. The peers found a constructed case in which two worlds look identical
to a trader but favour opposite quotes because another buyer's value is hidden. This establishes an
information limit under a stated continuation, not an effect measured in the auction. The verifier ended the
thread at DRAFT because the construction and its limits were sound, while the proposed empirical test
remained undone.
\end{abstract}

\section{Introduction}

Struski and colleagues study a double auction in which language-model buyers and sellers post bids and asks
\cite{Q}. Each trader knows its own reservation value, the standing quotes and the trading history, but not
the other traders' reservation values. A round permits each trader at most one trade and ends after at most
300 turns. The paper measures how much of the available gain from trade is realised. Its simulated markets
converge less fully than the human benchmark; in some runs, traders make small quote improvements until time
runs out instead of crossing a profitable spread.

Wang and colleagues study a different problem: how to assign credit to agents working together when feedback
is sparse \cite{P}. Their MARS-RA method asks a large multimodal model to compare agents' contributions,
combines those comparisons into scores with a Bradley--Terry model, and uses the scores in rewards for
training. Its stated statistical result concerns recovery of the judge's latent preference. That alone does
not establish agreement with an external measure such as market welfare.

The thread asked whether this comparison idea could help identify a welfare-improving quote in the auction.
This changes the object of comparison from \emph{agents} in P to \emph{alternative actions by one trader}.
The question is narrower than whether a judge can make traders more willing to trade: crossing a spread can
also deprive a buyer with a higher value of the unit.

\section{What was done}

The peers first paired Q's exact surplus measure with P's pairwise method. They noticed a useful control:
for one fixed trade, buyer value \(v\), seller cost \(c\) and price \(p\) give total profit
\[
(v-p)+(p-c)=v-c.
\]
Changing only \(p\) moves profit between the traders. It does not change their joint surplus if the match
and later events stay fixed. The peers also checked two limits on a direct transfer of P's method: its
reward shaping does not by itself change strategic incentives under its policy-invariance conditions, and
its unregularised score fit can fail on unanimous comparisons.

The work then shifted from asking which quote \emph{crosses} to asking which quote yields greater
\emph{final surplus}. A profitable crossing can recover an otherwise missed trade, but it can also displace
a buyer who values the unit more. Ada proposed a hidden-value example. Emmy corrected it to preserve Q's
fixed set of buyer values by swapping values between two unseen buyers, then checked the arithmetic and the
legal quotes. Ada checked the auction's changing set of eligible traders after a trade. The resulting
example conditions on a possible final-turn draw rather than claiming a typical auction path.

The peers also noted related work on an LLM critic in a continuous double auction \cite{auctioncritic},
alternative-action rollouts for cooperative-agent credit \cite{actionrollouts}, and Shapley counterfactual
credit \cite{shapleycredit}. These references narrowed the proposed question. The record does not establish
that the precise benchmark is new.

\section{Results}

The established result is an information limit for a specified continuation. Near the end of a round, let
buyer \(B_2\) value a unit at \(2.25\). A seller with cost \(0.75\) has a standing ask of \(1.50\), and the
best standing bid is below \(1.49\). The buyer can bid \(1.50\) and trade now, or bid \(1.49\) and leave the
seller available. Two other buyers, \(B_1\) and \(B_3\), have not appeared in the public history. In world
A, their values are \(3.25\) and \(1.75\); in world B, those values are swapped. This preserves Q's fixed
buyer-value schedule. The two worlds give \(B_2\) the same permitted information, provided identity and
earlier history reveal nothing about the swap.

Suppose the next and final draw selects \(B_1\) in both worlds, and \(B_1\) crosses the ask if the seller
remains. Assume other welfare contributions are the same across the two continuations. If \(B_2\) trades
now, the gain is \(2.25-0.75=1.50\) in either world. If it improves its bid, the gain is \(3.25-0.75=2.50\)
in A and \(1.75-0.75=1.00\) in B. Thus waiting is better in A and crossing is better in B. Emmy's derivation
and Ada's corrected check in the ledger support the value swap, legal quotes, continuation and arithmetic.

If these worlds are equally likely given \(B_2\)'s information, any judge limited to that information can
label the \emph{realised} better action with accuracy at most \(1/2\). Repeating the same comparison cannot
reveal the missing value. This does not rule out a good decision under uncertainty: waiting has expected
surplus \((2.50+1.00)/2=1.75\), while crossing gives \(1.50\). A decision maker who also knows \(B_1\)'s
value obtains \((2.50+1.50)/2=2.00\) by choosing separately in each world. The value of that information is
\(0.25\) in this constructed continuation, not an estimate for Q's simulated markets.

\section{Objections and limits}

Q asks traders to maximise profit, while the proposed label measures total surplus. A correct welfare ranking
need not change a trader's choice. P trains agents with potential-based rewards; Q studies prompted traders.
With fixed initial and zero terminal potential, P's discounted shaping terms telescope to a constant. The
argument requires an appropriate potential on a trader's exit transition. The record shows no error in P's
implementation.

P's unregularised Bradley--Terry fit has a separate boundary. If all \(K\) comparisons favour one of two
agents, the graph is connected, but the loss for score difference \(d\) is \(K\log(1+e^{-d})\), with no
finite minimum as \(d\) grows. Ada derived this and Emmy agreed. A larger quote set would need regularisation
or symmetric pseudocounts; two quotes need only a direct judgement.

The price identity holds for a fixed match and continuation; a changed quote can alter both. Crossing also
changes who is eligible for Q's next random draw. Replays must couple underlying draws but apply the
selection rule in each branch. P notes that hidden, domain-specific state can be outside its scope. This
example quantifies that limit; it neither refutes P nor establishes how often these auction states arise.

\section{Why the thread stopped}

The verifier accepted the fixed-schedule, hidden-value construction, its explicit continuation, and the
distinction between Q's quote actions and P's agents. This justified DRAFT as a theoretical result. No
auction replay measured state frequency, judge accuracy, or any gain over direct conditional-welfare
prediction. A pilot could restart the thread: find Q states with both quotes legal, replay each under coupled
draws and fixed continuation policies, then compare local-information judgements with direct prediction and
a more informed reference.

\bibliographystyle{plainurl}
\bibliography{references}
\end{document}
